\subsection{Семинар ``Unit-тестирование'' (2 очных часа)}

\textbf{Инструкция.}
Для своего варианта реализуйте указанную функцию и напишите набор
unit-тестов, проверяющих все существенные случаи её работы.

\subsubsection{Краткая теория}

\textbf{Unit-тест} --- автоматическая проверка отдельной функции на заранее
известных входных данных и ожидаемом результате.

В Python для написания unit-тестов можно использовать встроенный модуль
\texttt{unittest}.

\textbf{Структура теста}

Тесты обычно размещают в отдельном файле, имя которого начинается с
\texttt{test\_}:

\begin{verbatim}
from solution import classify_number
import unittest

class TestClassifyNumber(unittest.TestCase):

    def test_positive_even(self):
        self.assertEqual(
            classify_number(8),
            "positive even"
        )

if __name__ == "__main__":
    unittest.main()
\end{verbatim}

Каждый метод, имя которого начинается с \texttt{test\_}, является отдельным
тестом.

\textbf{Основные проверки}

Наиболее часто используются следующие методы:

\begin{itemize}
\item \texttt{assertEqual(actual, expected)} --- значения должны быть равны;
\item \texttt{assertNotEqual(actual, expected)} --- значения должны отличаться;
\item \texttt{assertTrue(value)} --- значение должно быть истинным;
\item \texttt{assertFalse(value)} --- значение должно быть ложным;
\item \texttt{assertIn(value, collection)} --- значение должно содержаться
в коллекции;
\item \texttt{assertRaises(Exception)} --- ожидается возникновение
указанного исключения. Используется как контекстный менеджер:
\begin{verbatim}
with self.assertRaises(ZeroDivisionError):
    divide(1, 0)
\end{verbatim}
\end{itemize}

Например:

\begin{verbatim}
self.assertEqual(is_perfect(6), True)
self.assertTrue(is_perfect(6))
self.assertFalse(is_perfect(10))
\end{verbatim}

\textbf{Что именно нужно проверять}

Хороший набор тестов должен проверять не только обычный случай. Для функции
следует определить различные классы входных данных и проверить каждый
существенный случай.

Особенно важны:

\begin{itemize}
\item обычные корректные значения;
\item граничные значения;
\item значения непосредственно до и после границы;
\item специальные случаи, например ноль или пустой список;
\item некорректные входные данные, если они предусмотрены заданием;
\item все ветви условных операторов \texttt{if}/\texttt{elif}/\texttt{else}.
\end{itemize}

Например, если функция использует условие

\begin{verbatim}
if x < 10:
...
elif x == 10:
...
else:
...
\end{verbatim}

тесты должны обеспечить выполнение каждой из трёх ветвей.

\textbf{Граничные значения}

При наличии диапазонов необходимо проверять сами границы и значения по обе
стороны от них.

Например, если скидка составляет 5\% при
$500 < price \leqslant 1000$, полезно проверить:

\begin{verbatim}
discount(500)
discount(501)
discount(1000)
discount(1001)
\end{verbatim}

Такие тесты позволяют обнаружить ошибки в операторах
\texttt{<}, \texttt{<=}, \texttt{>} и \texttt{>=}.

\textbf{Запуск тестов}

Отдельный тестовый файл можно запустить командой:

\begin{verbatim}
python -m unittest test_solution
\end{verbatim}

Для автоматического поиска всех тестов:

\begin{verbatim}
python -m unittest discover
\end{verbatim}

Если все тесты пройдены, программа сообщает об успешном выполнении.
Если тест не пройден, необходимо сравнить фактический результат функции
с ожидаемым и определить причину ошибки.


\begin{enumerate}

% 1
\item \texttt{classify\_triangle(a, b, c)}

Определяет тип треугольника по длинам трёх его сторон.

\begin{itemize}
    \item если треугольник невозможно построить, вернуть \texttt{"invalid"};
    \item если все три стороны равны, вернуть \texttt{"equilateral"};
    \item если ровно две стороны равны, вернуть \texttt{"isosceles"};
    \item если все три стороны различны, вернуть \texttt{"scalene"}.
\end{itemize}

Треугольник существует только при положительных длинах сторон и выполнении
строгого неравенства треугольника: каждая сторона должна быть строго меньше
суммы двух других. Вырожденный треугольник (где сумма двух сторон равна
третьей) считается несуществующим и также возвращает \texttt{"invalid"}.

Стороны задаются целыми числами, поэтому сравнение суммы двух сторон
с третьей выполняется точно, без погрешности.

% 2
\item \texttt{classify\_number(n)}

Определяет знак и чётность целого числа.

Функция должна возвращать одну из строк:
\begin{itemize}
    \item \texttt{"positive even"} --- положительное чётное число;
    \item \texttt{"positive odd"} --- положительное нечётное число;
    \item \texttt{"negative even"} --- отрицательное чётное число;
    \item \texttt{"negative odd"} --- отрицательное нечётное число;
    \item \texttt{"zero"} --- число равно нулю.
\end{itemize}

Отдельно проверьте положительные и отрицательные числа каждой чётности,
а также ноль.

% 3
\item \texttt{middle\_value(a, b, c)}

Возвращает среднее по величине из трёх чисел, то есть число, которое
после упорядочивания трёх значений занимает вторую позицию.

Функция не должна вычислять арифметическое среднее
$(a+b+c)/3$.

При наличии равных значений функция должна вернуть само значение,
занимающее среднюю позицию. Проверьте случаи с тремя равными числами,
с двумя равными числами и с тремя различными числами.

% 4
\item \texttt{max\_of\_three(a, b, c)}

Возвращает наибольшее из трёх чисел. Если несколько чисел равны между
собой и являются наибольшими, функция возвращает это значение.

Реализация должна использовать явные ветвления
\texttt{if}/\texttt{elif}/\texttt{else} без использования встроенных
функций (например, \texttt{max}).

Проверьте все варианты расположения наибольшего значения, а также случаи,
когда два или все три числа равны.

% 5
\item \texttt{is\_leap\_year(year)}

Определяет, является ли заданный год високосным.

Год является високосным, если:
\begin{itemize}
    \item он делится на 400; или
    \item он делится на 4, но не делится на 100.
\end{itemize}

Во всех остальных случаях год невисокосный.

Обязательно проверьте обычный високосный год, обычный невисокосный год,
год, кратный 100, но не кратный 400, и год, кратный 400.

% 6
\item \texttt{bmi\_category(weight, height)}

Вычисляет индекс массы тела:
\[
BMI = \frac{weight}{height^2},
\]
где \texttt{weight} задан в килограммах, а \texttt{height} --- в метрах.

Функция возвращает:
\begin{itemize}
    \item \texttt{"Underweight"} при $BMI < 18.5$;
    \item \texttt{"Normal"} при $18.5 \leq BMI < 25$;
    \item \texttt{"Overweight"} при $25 \leq BMI < 30$;
    \item \texttt{"Obese"} при $BMI \geq 30$.
\end{itemize}

Предполагается, что \texttt{weight} и \texttt{height} положительны.
Поведение при \texttt{height} $\leq 0$ или \texttt{weight} $\leq 0$
не определено и в тестах не проверяется.

Так как $BMI$ --- величина с плавающей точкой, при проверке границ
допускается сравнение с точностью до $10^{-9}$.

Проверьте значения непосредственно на каждой границе категорий и по обе
стороны от каждой границы.

% 7
\item \texttt{categorize\_temperature(temp)}

Определяет категорию температуры в градусах Цельсия. Функция работает
с непрерывными числовыми значениями:
\begin{itemize}
    \item \texttt{"freezing"} при $temp \leq 0$;
    \item \texttt{"cold"} при $0 < temp \leq 10$;
    \item \texttt{"cool"} при $10 < temp \leq 20$;
    \item \texttt{"warm"} при $20 < temp \leq 30$;
    \item \texttt{"hot"} при $temp > 30$.
\end{itemize}

Проверьте значения на границах и непосредственно по обе стороны от них.

% 8
\item \texttt{triangle\_angle\_type(a, b, c)}

Определяет вид треугольника по длинам его сторон через соотношение
между квадратами сторон.

Функция возвращает:
\begin{itemize}
    \item \texttt{"invalid"}, если треугольник не существует;
    \item \texttt{"right"}, если треугольник прямоугольный;
    \item \texttt{"acute"}, если все его углы острые;
    \item \texttt{"obtuse"}, если один из углов тупой.
\end{itemize}

Для существующего треугольника сначала определите наибольшую сторону
$c$. Затем сравнивайте $c^2$ с $a^2+b^2$.

Стороны задаются целыми числами, поэтому все сравнения выполняются точно.

Проверьте также вырожденные и другие невозможные тройки сторон.

% 9
\item \texttt{quadrant(x, y)}

Определяет положение точки $(x,y)$ относительно координатных осей.

Функция возвращает:
\begin{itemize}
    \item \texttt{1} --- точка в первой четверти;
    \item \texttt{2} --- точка во второй четверти;
    \item \texttt{3} --- точка в третьей четверти;
    \item \texttt{4} --- точка в четвёртой четверти;
    \item \texttt{"origin"} --- точка $(0,0)$;
    \item \texttt{"axis"} --- точка лежит на одной из координатных осей,
    но не является началом координат.
\end{itemize}

Проверьте точки во всех четырёх четвертях, на обеих осях и в начале
координат.

% 10
\item \texttt{days\_in\_month(month, leap)}

Возвращает количество дней в указанном месяце.

\texttt{month} --- номер месяца от 1 до 12.
\texttt{leap} --- признак високосного года.

Количество дней:
\begin{itemize}
    \item 31 день --- январь, март, май, июль, август, октябрь, декабрь;
    \item 30 дней --- апрель, июнь, сентябрь, ноябрь;
    \item февраль --- 29 дней при \texttt{leap=True} и 28 дней при
    \texttt{leap=False}.
\end{itemize}

Гарантируется, что \texttt{month} находится в диапазоне от 1 до 12.
Поведение при других значениях \texttt{month} не определено и в тестах
не проверяется.

Проверьте каждый тип месяца и оба варианта февраля.

% 11
\item \texttt{traffic\_fine(speed, zone)}

Вычисляет размер штрафа за превышение скорости.

Скорость в \texttt{speed} задаётся в км/ч.
Для каждой зоны установлено следующее ограничение скорости:
\begin{itemize}
    \item \texttt{"residential"} --- 40 км/ч;
    \item \texttt{"city"} --- 60 км/ч;
    \item \texttt{"highway"} --- 110 км/ч.
\end{itemize}

Сначала вычисляется превышение:
\[
excess = speed - speed\_limit.
\]

Размер штрафа:
\begin{itemize}
    \item \texttt{"residential"}:
    \begin{itemize}
        \item $excess \leq 10$ --- 0;
        \item $10 < excess \leq 20$ --- 100;
        \item $excess > 20$ --- 200;
    \end{itemize}

    \item \texttt{"city"}:
    \begin{itemize}
        \item $excess \leq 15$ --- 0;
        \item $15 < excess \leq 30$ --- 75;
        \item $excess > 30$ --- 150;
    \end{itemize}

    \item \texttt{"highway"}:
    \begin{itemize}
        \item $excess \leq 20$ --- 0;
        \item $20 < excess \leq 40$ --- 50;
        \item $excess > 40$ --- 100.
    \end{itemize}
\end{itemize}

Если \texttt{speed} меньше или равна ограничению, превышение
отрицательно или равно нулю, и штраф равен $0$.
Если \texttt{speed} отрицательна, превышения нет, штраф равен $0$.

Если зона не равна одной из трёх допустимых строк, функция возвращает
\texttt{"invalid zone"}.

Проверьте отсутствие превышения, все границы штрафов и значения сразу
по обе стороны каждой границы для каждой зоны.

% 12
\item \texttt{compare\_three\_numbers(a, b, c)}

Сравнивает три числа и возвращает:
\begin{itemize}
    \item \texttt{"all equal"} --- все три значения равны;
    \item \texttt{"two equal"} --- ровно два значения равны;
    \item \texttt{"all different"} --- все три значения различны.
\end{itemize}

Проверьте все три возможных случая независимо от того, какая именно
пара значений совпадает.

% 13
\item \texttt{max\_digit(n)}

Возвращает наибольшую цифру целого числа.

Знак числа не учитывается: для отрицательного числа рассматривается
абсолютное значение.

Например, для $n=-582$ результатом является $8$.
Для числа $0$ результат равен $0$.

Проверьте однозначные числа, многозначные числа, повторяющиеся цифры,
а также положительные и отрицательные значения.

% 14
\item \texttt{triangle\_area(a, b, c)}

Вычисляет площадь треугольника по длинам его сторон, используя формулу
Герона.

Если треугольник не существует (включая вырожденный случай, когда сумма
двух сторон равна третьей), функция возвращает \texttt{"invalid"}.

Стороны задаются целыми числами. Площадь возвращается как число
с плавающей точкой; при проверке в тестах допускается сравнение
с точностью до $10^{-6}$.

Проверьте корректные треугольники, вырожденные случаи и заведомо
невозможные тройки сторон.

% 15
\item \texttt{next\_day(day, month, leap)}

Возвращает дату следующего календарного дня.

\texttt{day} --- номер дня месяца, \texttt{month} --- номер месяца,
\texttt{leap} --- признак високосного года.

В обычный день увеличивается только номер дня.
Если текущий день является последним днём месяца, результатом является
первый день следующего месяца.
После 31 декабря результатом является 1 января.

Количество дней в феврале зависит от значения \texttt{leap}.

Гарантируется, что входная дата корректна: \texttt{month} от 1 до 12,
\texttt{day} от 1 до количества дней в этом месяце с учётом \texttt{leap}.

Проверьте обычный переход внутри месяца, переход на следующий месяц,
переход с февраля, переход с 31 декабря и оба варианта високосности
февраля.

% 16
\item \texttt{is\_inside\_rectangle(x, y, x1, y1, x2, y2)}

Проверяет, принадлежит ли точка $(x,y)$ прямоугольнику, заданному двумя
противоположными вершинами $(x1,y1)$ и $(x2,y2)$.

Граница прямоугольника считается частью прямоугольника.

Функция возвращает \texttt{True}, если точка находится внутри
прямоугольника или на его границе, и \texttt{False} в противном случае.

Координаты противоположных вершин могут быть заданы в любом порядке.

% 17
\item \texttt{discount(price)}

Вычисляет итоговую цену товара после применения скидки.

Правила:
\begin{itemize}
    \item при $price > 1000$ предоставляется скидка 10\%;
    \item при $500 < price \leq 1000$ предоставляется скидка 5\%;
    \item при $price \leq 500$ скидка не предоставляется.
\end{itemize}

Функция возвращает именно цену после скидки, а не размер скидки.

Предполагается, что \texttt{price} неотрицательна. Поведение при
отрицательной цене не определено и в тестах не проверяется.

Так как \texttt{price} может быть числом с плавающей точкой,
при проверке результата допускается сравнение с точностью до $10^{-9}$.

Проверьте значения на обеих границах и непосредственно по обе стороны
от них.

% 18
\item \texttt{closest\_to\_zero(lst)}

Возвращает элемент списка, абсолютное значение которого минимально.

Если два или более элемента находятся на одинаковом расстоянии от нуля,
возвращается первый из них в исходном списке.

Список содержит хотя бы один элемент.

Проверьте положительные и отрицательные значения, наличие нуля,
а также случай равного расстояния до нуля.

% 19
\item \texttt{season(month)}

По номеру месяца возвращает название соответствующего времени года:
\begin{itemize}
    \item \texttt{"Winter"} --- декабрь, январь, февраль;
    \item \texttt{"Spring"} --- март, апрель, май;
    \item \texttt{"Summer"} --- июнь, июль, август;
    \item \texttt{"Autumn"} --- сентябрь, октябрь, ноябрь.
\end{itemize}

Гарантируется, что \texttt{month} находится в диапазоне от 1 до 12.
Поведение при других значениях не определено и в тестах не проверяется.

Проверьте все месяцы, включая месяцы, являющиеся границами сезонов.

% 20
\item \texttt{simple\_calculator(a, b, op)}

Выполняет одну арифметическую операцию над числами $a$ и $b$.

Параметр \texttt{op} определяет операцию:
\begin{itemize}
    \item \texttt{+} --- $a+b$;
    \item \texttt{-} --- $a-b$;
    \item \texttt{*} --- $a\cdot b$;
    \item \texttt{/} --- $a/b$.
\end{itemize}

При попытке деления на ноль функция возвращает \texttt{"error"}.
При неизвестном операторе функция также возвращает \texttt{"error"}.

Проверьте каждую операцию, деление на ноль и недопустимый оператор.

% 21
\item \texttt{sum\_positive(lst)}

Возвращает сумму всех строго положительных элементов списка.

Ноль и отрицательные элементы не учитываются.

Если в списке нет положительных элементов, результатом является $0$.

Проверьте список с положительными и отрицательными числами, список без
положительных чисел и список, содержащий нули.

% 22
\item \texttt{sum\_even(lst)}

Возвращает сумму всех чётных элементов списка.

Чётность определяется для целых чисел; отрицательные чётные числа также
учитываются.

Если чётных элементов нет, функция возвращает $0$.

Проверьте положительные, отрицательные, нулевые и нечётные элементы.

% 23
\item \texttt{sum\_odd(lst)}

Возвращает сумму всех нечётных элементов списка.

В сумму включаются как положительные, так и отрицательные нечётные числа.
Ноль не является нечётным и в сумму не включается.

Если нечётных элементов нет, функция возвращает $0$.

Проверьте смешанный список, список без нечётных чисел и наличие
отрицательных значений.

% 24
\item \texttt{reverse\_signs(lst)}

Возвращает новый список, в котором каждый элемент исходного списка
заменён на противоположный по знаку:
\[
x \mapsto -x.
\]

Ноль остаётся равным нулю.

Исходный список не должен изменяться.

Проверьте положительные, отрицательные и нулевые элементы, а также
пустой список.

% 25
\item \texttt{nearest\_multiple(n, m)}

Возвращает целое число, кратное $m$, которое находится на минимальном
расстоянии от $n$.

Параметры: $n$ --- целое число, $m$ --- целое положительное число
($m>0$). Множество возможных результатов имеет вид $k\cdot m$,
где $k$ --- целое число.

Если два кратных находятся от $n$ на одинаковом расстоянии, возвращается
меньшее из них.

Проверьте случаи, когда $n$ уже кратно $m$, находится между двумя
кратными, а также случай равного расстояния.

% 26
\item \texttt{sort\_three(a, b, c)}

Возвращает тройку чисел в порядке неубывания:
от наименьшего значения к наибольшему.

Например, для входа $(3,1,2)$ результатом является $(1,2,3)$.

Реализация должна использовать ветвления
\texttt{if}/\texttt{elif}/\texttt{else}.

Повторяющиеся значения должны сохраняться в результате.
Проверьте все варианты порядка трёх различных чисел и случаи
с повторениями.

% 27
\item \texttt{validate\_password(password)}

Проверяет, соответствует ли строка пароля всем трём требованиям:
\begin{itemize}
    \item длина строки не менее 8 символов;
    \item строка содержит хотя бы одну цифру от \texttt{0} до \texttt{9};
    \item строка содержит хотя бы одну заглавную латинскую букву
    от \texttt{A} до \texttt{Z}.
\end{itemize}

Функция возвращает \texttt{True}, если выполнены все три условия,
и \texttt{False} в противном случае.

Проверьте корректный пароль, а также случаи нарушения каждого требования
по отдельности и нескольких требований одновременно.

% 28
\item \texttt{is\_perfect(n)}

Проверяет, является ли целое положительное число совершенным.

Число называется совершенным, если оно равно сумме всех своих
положительных собственных делителей, то есть делителей, меньших самого
числа.

Например, $6=1+2+3$, поэтому 6 является совершенным числом.

Для $n\leq0$ функция возвращает \texttt{False}.

Проверьте известное совершенное число, обычное несовершенное число,
простое число и значения, меньшие или равные нулю.

% 29
\item \texttt{sum\_digits(n)}

Возвращает сумму всех цифр целого числа.

Знак числа не учитывается: для отрицательного числа рассматривается
абсолютное значение.

Например, для $-305$ результатом является $3+0+5=8$.

Для $0$ результат равен $0$.

Проверьте однозначные и многозначные числа, наличие нулей внутри числа
и отрицательные значения.

% 30
\item \texttt{count\_vowels(s)}

Возвращает количество английских гласных букв в строке.

Гласными считаются символы:
\texttt{a}, \texttt{e}, \texttt{i}, \texttt{o}, \texttt{u}
в любом регистре.

Все остальные символы, включая согласные, цифры, пробелы и знаки
препинания, не учитываются.

Проверьте строку без гласных, строку только из гласных, смешанный текст,
разный регистр и пустую строку.

% 31
\item \texttt{is\_palindrome(s)}

Проверяет, является ли строка палиндромом.

Палиндромом считается строка, которая совпадает со своей записью
в обратном порядке.

Сравнение выполняется посимвольно с учётом регистра и пробелов.
Например, \texttt{"level"} является палиндромом, а
\texttt{"Level"} --- нет.

Пустая строка и строка из одного символа считаются палиндромами.

Проверьте палиндромы чётной и нечётной длины, непалиндромные строки,
строку из одного символа и пустую строку.

% 32
\item \texttt{remove\_duplicates(lst)}

Возвращает новый список, содержащий каждый встречающийся элемент
ровно один раз.

Порядок элементов должен сохраняться: при повторном появлении элемента
остаётся его первое вхождение.

Например,
\[
[3,1,3,2,1] \rightarrow [3,1,2].
\]

Исходный список не должен изменяться.

Проверьте список без повторов, список со множественными повторами,
список из одинаковых элементов и пустой список.

% 33
\item \texttt{factorial\_iterative(n)}

Вычисляет факториал неотрицательного целого числа:
\[
n! = 1\cdot2\cdot3\cdots n.
\]

По определению,
\[
0! = 1.
\]

Функция должна использовать цикл и не должна использовать рекурсию.

Параметр $n$ --- целое число. Для отрицательного $n$ функция должна
возвращать \texttt{"error"}.

Проверьте $0$, $1$, несколько обычных значений и отрицательное число.

% 34
\item \texttt{fizz\_buzz(n)}

Возвращает список результатов для всех целых чисел от 1 до $n$ включительно.

Правила замены:
\begin{itemize}
    \item число, кратное 3, заменяется на \texttt{"Fizz"};
    \item число, кратное 5, заменяется на \texttt{"Buzz"};
    \item число, кратное одновременно 3 и 5, заменяется на
    \texttt{"FizzBuzz"};
    \item остальные числа остаются числами.
\end{itemize}

Проверка кратности 15 должна иметь приоритет над отдельной проверкой
кратности 3 или 5.

Для $n\leq0$ функция возвращает пустой список.

Проверьте значения до первой замены, значения, кратные только 3,
только 5 и одновременно 3 и 5, а также границу диапазона.

% 35
\item \texttt{is\_prime(n)}

Проверяет, является ли целое число простым.

Простым считается натуральное число, имеющее ровно два положительных
делителя: 1 и само число.

Числа меньше 2 простыми не являются.

Параметр $n$ --- целое число.

Функция возвращает \texttt{True} для простого числа и \texttt{False}
для составного или недопустимого числа.

Проверьте числа меньше 2, число 2, нечётные простые числа, чётные
составные числа, составные нечётные числа и несколько чисел с делителями.

\end{enumerate}